%%
%% spring-lecture02.tex
%% 
%% Made by Alex Nelson <pqnelson@gmail.com>
%% Login   <alex@lisp>
%% 
%% Started on  2026-04-09T09:56:56-0700
%% Last update 2026-04-09T09:56:56-0700
%% 

\lecture[Smooth functions and differentials]{}

\begin{definition}[Smooth function]
Let $M$ and $N$ be smooth manifolds, let $f\colon M\to N$ be any
continuous function. We say $f$ is \define{Smooth} if for any local
charts $(U,\varphi)\in\smoothstr{M}$ and $(V,\psi)\in\smoothstr{N}$
such that for any $x\in U$ and $f(x)\in V$ and $f(U)\subset V$ we have
\begin{equation}
\psi\circ f\circ\varphi^{-1}\colon\varphi(U)\to\psi(f(V))
\end{equation}
is smooth (in the sense of calculus).
\end{definition}

\begin{node}
We n eed a way to infinitesimally perturb a point of a manifold in
order to determine the first variation for a function. For calculus,
this is precisely the derivative
\begin{equation}
\frac{\D f}{\D x}(x)=\lim_{h\to 0}\frac{f(x+h)-f(x)}{h},
\end{equation}
but what does $f(x+h)-f(x)$ even mean for functions on manifolds? Or
dividing by $h$? What even is $h$?
\end{node}

\begin{definition}[Tangent vectors]
Let $M$ be a smooth $n$-manifold, let $x\in M$. We define the
\define{Set of Curves Passing Through $x$}
\begin{equation}
\mathcal{C}(M,x)=\{\gamma\colon I\to M\mid 0\in I, x=\gamma(0)\}
\end{equation}
where $I$ is an open interval $(a,b)\subset\RR$.

On $\mathcal{C}(M,x)$ we can define the following equivalence
relation: $\gamma_{1}\sim\gamma_{2}$ if there exists some local chart
$(U,\varphi)$ containing $x\in U$ such that
$(\varphi\circ\gamma_{1})'(0)=(\varphi\circ\gamma_{2})'(0)$ their
velocity vectors at $x$ are the same in the local chart.

Using these notions, we define the \define{Tangent Space} of $M$ at
$x$ to be
\begin{equation}
T_{x}M:=\{[\gamma]\mid \gamma\in\mathcal{C}(M,x)\}
\end{equation}
the collection of equivalence classes of $\mathcal{C}(M,x)$ under the
previous paragraph's equivalence relation. We call the elements of
$T_{x}M$ \define{Tangent Vectors}.
\end{definition}

\begin{xca}
Prove this equivalence relation $\gamma_{1}\sim\gamma_{2}$ is
independent of choice of chart $(U,\varphi)$.
\end{xca}

\begin{node}
We want to think of $T_{x}M$ as the set of velocity vectors of curves
passing through $x$, but wherre do the velocity vectors live? If we
have a submanifold of $\RR^{N}$, then they live in the ``obvious''
tangent vector space. Here, we \emph{borrow} the vector space from the
local charts.
\end{node}

\begin{remark}
When we have a local chart $(U,\varphi)$ at a point $x\in M$, there
are preferred tangent vectors. If $M$ is a smooth $n$-manifold, then
we may write the components of $\varphi$ as
\begin{equation}
\varphi=(\xi^{1},\dots,\xi^{n}),
\end{equation}
where each $\xi^{i}$ is a smooth function $\xi^{i}\colon U\to\RR$.
Then we define the \define{$i^{\text{th}}$ Coordinate Tangent Vector}
at $x$ in the local chart $(U,\varphi)$ to be
\begin{equation}
\frac{\partial}{\partial\xi^{i}}(x)=[t\mapsto\varphi^{-1}(\varphi(x)+te_{i})]
\end{equation}
the equivalenc e class of curves $t\mapsto\varphi^{-1}(\varphi(x)+te_{i})$
where $e_{i}\in\RR^{n}$ is the $i^{\text{th}}$ standard basis vector
of $\RR^{n}$.
\end{remark}

\begin{definition}[Differential]
Let $M$ and $N$ be smooth manifolds. Let $f\colon M\to N$ be a smooth map.
Given a point $x\in M$, the \define{Differential} of $f$ at $x$ is the
linear map
\begin{equation}
\D f(x)\colon T_{x}M\to T_{f(x)}N
\end{equation}
given by
\begin{equation}
\D f(x)[v]=[f\circ\gamma]
\end{equation}
where $v=[\gamma]$. Note the square brackets on the left-hand side
mean $\D f(x)$ is a linear operator, but the square brackets on the
right-hand side refer to the equivalence class of curves.

Alternate notation you may find: $(\D f)_{x}(v)=[f\circ\gamma]$.
\end{definition}

\begin{node}
Why is this useful/ We can extend all the results from calculus to
manifolds. For example, the inverse function theorem (if $\D f$ is
invertible at $x$, then there is a neighborhood where $f$ is
invertible), etc.
\end{node}

\begin{definition}[Immersion, submersion, local diffeomorphism, embeddings]
Let $M$ and $N$ be smooth manifolds. Let $f\colon M\to N$ be a smooth map.
Let $x\in M$ be any point.
\begin{enumerate}
\item We say $f$ is a \define{Immersion} at $x$ if and only if $\D f(x)$
is injective; we say $f$ is an ``immersion'' if it is an immersion at
every point of $M$
\item We say $f$ is a \define{Submersion} at $x$ if and only if $\D f(x)$
is surjective; we say $f$ is a ``submersion'' if it is a submersion at
every point of $M$
\item We say $f$ is a \define{Local Diffeomorphism} at $x$ if and only
  if $\D f(x)$ is invertible. (This secrtetly contains the inverse
  function theorem)
\item We say $f$ is an \define{Embedding} if and only if $f$ is an
  immersion (at every point) and $f$ is a homeomorphism onto its image.
\end{enumerate}
Note: ``$f$ is a homeomorphism onto its image'' is a global condition.
\end{definition}

\begin{definition}[Diffeomorphism]
Let $f\colon M\to N$ be a function between two smooth manifolds. We
say $f$ is a \define{Diffeomorphism} if $f$ is smooth, invertible, and
has a smooth inverse. We write $M\iso N$ to indicate $M$ is
diffeomorphic to $N$ (i.e., there exists a diffeomorphism $f\colon M\to N$)
and it forms an equivalence relation.
\end{definition}

\begin{proposition}
Let $M$ be a smooth $n$-manifold. Let $f\colon M\to\RR^{k}$ be an
embedding. Then $f(M)\subset\RR^{k}$ is a smooth $n$-submanifold of
$\RR^{k}$ and $M\iso f(M)$.
\end{proposition}

\begin{theorem}[Whitney's embedding theorem]
Every smooth $n$-manifold is an $n$-submanifold of $\RR^{2n}$.
\end{theorem}

\begin{node}
If you want a smooth immersion, then you can make $M$ an
$n$-submanifold of $\RR^{2n-1}$.
\end{node}